\documentclass[10pt,a4paper]{amsart}
\usepackage{amsmath,amssymb,amsthm,graphicx}
\usepackage[scr=rsfs,cal=euler]{mathalfa}
\multlinegap0pt
\newcommand{\RedefinitSymbole}[1]{%
\expandafter\let\csname old\string#1\endcsname=#1
\let#1=\relax
\newcommand{#1}{\csname old\string#1\endcsname\,}%
}
\RedefinitSymbole{\forall} \RedefinitSymbole{\exists}

\newcommand{\llbracket}{[\![}
\newcommand{\rrbracket}{]\!]}

\newcommand{\spfrac}[2]{\sfrac{#1}{(#2)}}

\newcommand\mto{\mathchoice{\longmapsto}{\mapsto}{\mapsto}{\mapsto}}

\newcommand{\bn}{\boldsymbol{n}}
\newcommand{\bT}{\boldsymbol{T}}
\newcommand{\bV}{\boldsymbol{V}}
\newcommand{\loc}{\mathrm{loc}}
\newcommand{\nc}{\mathrm{nc}}
\newcommand{\obs}{\mathrm{obs}}

\DeclareMathOperator{\Id}{Id}
\DeclareMathOperator{\Ker}{Ker}
\DeclareMathOperator{\Ran}{Ran}
\DeclareMathOperator{\Span}{Span}
\usepackage{bbm}
\usepackage{tikz}
\usetikzlibrary{hobby}

\newcommand\R{{\mathbb R}}
\newcommand\N{{\mathbb N}}

\newcommand{\Hreg}{\mathrm{H}_{\omega,\Theta}}
\newcommand{\norm}[1]{\Vert#1\Vert}
\newcommand{\bnorm}[1]{\bigl \Vert#1\bigr\Vert}
\newcommand{\Bnorm}[1]{\Bigl\Vert#1\Bigr\Vert}
\newcommand{\bbnorm}[1]{\biggl\Vert#1\biggr\Vert}

\newtheorem{theorem}{Theorem}[section]
\newtheorem{corollary}[theorem]{Corollary}
\newtheorem{proposition}[theorem]{Proposition}
\newtheorem{lemma}[theorem]{Lemma}
\theoremstyle{definition}
\newtheorem{definition}[theorem]{Definition}
\newtheorem{remark}[theorem]{Remark}
\title{Exact heat-shape desensitization when the control contains the observation boundary}
\author{Sylvain Ervedoza, Pierre Lissy, Yannick Privat}
\date{doi:10.5802/jep.209; CC BY 4.0}
\begin{document}
\maketitle
\noindent\textit{Editorial note.} Original author statement and full paired proof retained; source license and version reused. No mathematical modifications or added assumptions.\par\medskip
\noindent\textit{Editorial note.} Public 2109 is finite-dimensional exactness plus global approximation in disjoint geometry. The present target is all-direction exactness in a boundary-collar geometry, obtained by a separate cutoff/commutator construction; the interior-observation easy case is not counted.\par\medskip
\noindent\textit{Allowed standard prerequisites (editorial).}\begin{itemize}
\item Classical heat equation and functional analysis. Weak and transposition solutions, Sobolev trace maps and parabolic regularity, Holmgren unique continuation, Hilbert-space range/kernel duality, standard heat observability and convex minimization. The displayed Hadamard derivative formula is a standard supplied input. Target-specific constructions remain in the original proof.
\end{itemize}\medskip
Let $T>0$ denote a horizon of time, $\omega$ and $\Theta$ be two open subsets of $\R^d$, $d\in\mathbb N^*$, and $\xi \in L^2(0,T; L^2(\R^d))$.

For a bounded, connected and open set $\Omega$ of $\mathbb{R}^d$ of class $\mathscr{C}^2$, we consider the shape functional $J_h$, indexed by $h \in L^2(0,T; L^2(\omega))$, defined by
\begin{equation}
\label{def-J}
J_h(\Omega) = \dfrac{1}{2}\int_{0}^{T} \int_{\Theta}y_{\Omega,h}(t,x)^{2}\, dx\,dt,
\end{equation}
where $y_{\Omega,h}\in L^2((0,T)\times \mathbb R^d)$ is defined on $(0,T)\times\Omega$ as the unique weak solution of
\begin{equation}
\label{y-Omega-h}
\left\{
\begin{array}{rcll}
\dfrac{\partial y}{\partial t}-\Delta y &=&\xi +h\mathbbm{1}_{\omega}\quad&\text{ in }(0,T) \times \Omega,
\\
y&=& 0 &\text{ on }(0,T) \times \partial \Omega,
\\
y\left(0,\cdot\right) &=& 0 &\text{ in }\Omega,
\end{array}
\right.
\end{equation}
($\mathbbm{1}_\omega$ denotes the characteristic function of the set $\omega$) extended by $0$ outside $(0,T) \times \Omega$.

In the sequel, given $\bV\in W^{3,\infty}(\R^d, \R^d)$, we introduce the family $\{\Omega_{\tau\bV}\}_{\tau\in [0,\tau_{\bV})}$ of domains defined by
\[
\Omega_{\tau\bV}= (\Id+\tau\bV)(\Omega_0).
\]

\begin{proposition}[{\cite[Proof of Prop.\,1.1]{Lissy-Privat-Simpore}}]
\label{Prop-J-differentiable}
Let $\xi \in L^2(0,T; L^2(\R^d))$ and $h \in L^2(0,T; L^2(\omega))$. For all $ \bV \in W^{3, \infty} (\R^d, \R^d)$, the mapping $\tau \mto J_h(\Omega_{\tau\bV})$ is differentiable at $\tau = 0$ and
\begin{equation}
\dfrac{d }{d \tau}\left({J}_h\left(\Omega_{\tau\bV} \right)\right)\Big|_{\tau=0}
\label{Diff-J-V-2}
=
\int_{\partial\Omega_0} \bV \cdot \bn \biggl(\int_0^T \partial_n y_0 \partial_n z_0 \, dt\biggr) d\sigma,
\end{equation}
where the pair $(y_0,z_0)$ solves the coupled forward-backward system
\begin{equation}
\label{y-Omega-0-h}
\left\{
\begin{array}{rcll}
\dfrac{\partial y_0}{\partial t}-\Delta y_0 &=&\xi +h\mathbbm{1}_{\omega}\quad&\text{ in }(0,T) \times \Omega_0,
\\
y_0 &=& 0 &\text{ on }(0,T) \times \partial \Omega_0,
\\
y_0\left(0,\cdot\right) &=& 0 &\text{ in }\Omega_0,
\\
-\dfrac{\partial z_0}{\partial t}-\Delta z_0 &=&\mathbbm{1}_{\Theta}y_0 &\text{ in }(0,T) \times \Omega_0,
\\
z_0 &=& 0 &\text{ on }(0,T) \times \partial \Omega_0,
\\
z_0\left(T,\cdot\right) &=& 0 &\text{ in }\Omega_0.
\end{array}
\right.
\end{equation}
($\mathbbm{1}_\Theta$ denotes the characteristic function of the set $\Theta$.)
\end{proposition}

\begin{definition}
Let {$\xi \in L^2(0,T; L^2(\R^d))$}.
We say that the control function $h \in L^2(0,T; L^2(\omega))$ desensitizes $J_h$ exactly at $\Omega_0$ at the first order with respect to domain variations if
\begin{equation}
\label{Exact-Insentitizing}
\forall\bV\in W^{3,\infty}(\R^d, \R^d),
\quad
\dfrac{d }{d \tau}\left(J_h \left(\Omega_{\tau\bV} \right)\right)\Big|_{\tau=0}=0.
\end{equation}

Let $\mathcal{E}$ be a linear subspace of $W^{3,\infty}(\R^d, \R^d)$.
We say that the control function $h \in L^2(0,T; L^2(\omega))$ exactly desensitizes $J_h$ for $\mathcal{E}$ at $\Omega_0$ at the first order with respect to domain variations if
\begin{equation}
\label{Finite-Insensitizing}
\forall\bV\in \mathcal{E},
\quad
\dfrac{d }{d \tau}\left(J_h \left(\Omega_{\tau\bV} \right)\right)\Big|_{\tau=0}=0.
\end{equation}

Given $\varepsilon>0$, we say that the control function $h \in L^2(0,T; L^2(\omega))$ $\varepsilon$-approximately desensitizes $J_h$ at $\Omega_0$ at the first order with respect to domain variations if
\begin{equation}
\label{Approx-Insensitizing}
\forall\bV\in W^{3,\infty}(\R^d, \R^d),
\quad
\left|\dfrac{d }{d \tau}\left(J_h\left(\Omega_{\tau \bV}\right)\right)\Big|_{\tau=0}\right |\leq \varepsilon \| \bV \|_{W^{3, \infty}(\R^d; \R^d)}.
\end{equation}
\end{definition}

\begin{theorem}
\label{Thm-insensitizing-Boundary}
Let $\omega$ and $\Theta$ be smooth non-empty open subsets of $\Omega_0$ such that
\begin{equation}
\label{Assumption-insensitizing-Boundary}
\partial \Theta \text{ has only one connected component, } \quad \Theta \Subset \Omega_0, \quad \text{ and } \quad \partial \Theta \subset \omega.
\end{equation}
Then, for all $\xi \in L^2(0,T; L^2(\Omega_0))$, there exists $h \in L^2(0,T; L^2(\omega))$ such that the functional $J_h$ in \eqref{def-J} is exactly desensitized in the sense of \eqref{Exact-Insentitizing}.
\end{theorem}

\subsection{Proof of Theorem \ref{Thm-insensitizing-Boundary}}
\label{Sec-Proof-Thm-insensitizing-Boundary}

We start by introducing open sets $\omega_0$, $\omega_1$, $\omega_2$, and~$\omega_3$ such that
\begin{equation*}
\partial \Theta \subset \omega_0 \Subset \omega_1 \Subset \omega_2 \Subset \omega_3 \Subset \omega,
\end{equation*}
which is possible thanks to Assumption \eqref{Assumption-insensitizing-Boundary}.
We also introduce a smooth function $\eta_{23} = \eta_{23}(x)$ taking value $1$ in $\Omega \setminus \overline{\omega_3}$ and vanishing in $\overline{\omega_2}$.

We fix $\xi \in L^2(0,T; L^2(\R^d))$, and we introduce the solution $y_\xi$ to
\begin{equation*}
\left\{
\begin{array}{rcll}
\partial_t y_\xi - \Delta y_\xi &=& \eta_{23} \xi, \quad & (t,x) \in (0,T) \times \Omega_0,
\\
y_\xi(t, x) &=& 0, \quad & (t,x) \in (0,T) \times \partial \Omega_0,
\\
y_\xi(0,x) &=& 0, \quad & x \in \Omega,
\end{array}
\right.
\end{equation*}
and the solution $z_\xi$ to
\begin{equation}
\label{Heat-Eq-Adjoint-z-f}
\left\{
\begin{array}{rcll}
- \partial_t z_\xi - \Delta z_\xi & =& \eta_{12} y_\xi \mathbbm{1}_{\Theta}, \quad & (t,x) \in (0,T) \times \Omega_0,
\\
z_\xi(t, x) &=& 0, \quad & (t,x) \in (0,T) \times \partial \Omega_0,
\\
z_\xi(T,x) &=& 0, \quad & x \in \Omega_0.
\end{array}
\right.
\end{equation}
where $\eta_{12} \!=\! \eta_{12}(x)$ is a smooth function taking value $1$ in $\Omega_0 \setminus \overline{\omega_2}$ and vanishing in~$\overline{\omega_1}$.

Then, we introduce a smooth function $\eta_{01} = \eta_{01}(x)$ such that $\eta_{01}$ vanishes in $ \omega_0$ and equal to $1$ in $\Omega_0 \setminus \omega_1$, so that $\mathbbm{1}_\Theta \eta_{01}$ is actually a smooth function taking value $1$ in $\Theta \setminus \omega_1$ and vanishing in $\omega_0 \cup (\Omega_0 \setminus \Theta)$. Then $ z_0(t,x)= \mathbbm{1}_\Theta(x) \eta_{01}(x) z_\xi(t,x)$ satisfies
\begin{equation*}
\left\{
\begin{array}{rcll}
- \partial_t z_0 - \Delta z_0 &=& \eta_{01} \eta_{12} y_\xi \mathbbm{1}_{\Theta}
- [\Delta, \mathbbm{1}_\Theta \eta_{01}] z_\xi
, \quad & (t,x) \in (0,T) \times \Omega_0,
\\
z_0(t, x) &=& 0, \quad & (t,x) \in (0,T) \times \partial \Omega_0,
\\
z_0(T,x) &=& 0, \quad & x \in \Omega_0,
\end{array}
\right.
\end{equation*}
and, by construction,
\[
z_0(t,x) = 0 \quad\text{for all } (t,x) \in (0,T) \times (\Omega_0 \setminus \Theta).
\]
Now, we remark that by construction, $[\Delta, \mathbbm{1}_\Theta \eta_{01}] z_\xi$ is localized in $\omega_1 \cap \Theta$. Besides, local regularity results for \eqref{Heat-Eq-Adjoint-z-f} imply that $z_\xi$ is $\mathscr{C}^2([0,T] \times \overline{\omega_1 \cap \Theta})$. We then take
\[
y_0(t,x) = \eta_{01}(x) \eta_{12}(x) y_\xi(t,x) - [\Delta, \mathbbm{1}_\Theta \eta_{01} ] z_\xi(t,x), \quad \text{for all } (t,x) \in (0,T) \times \Omega_0,
\]
which satisfies \eqref{y-Omega-0-h} for $h$ given for $(t,x) \in (0,T) \times \Omega_0$ by
\begin{multline*}
h(t,x) = (\eta_{01}(x) \eta_{12}(x) \eta_{23}(x)-1) \xi(t,x) - [\Delta, \eta_{01} \eta_{12}] y_\xi(t,x)\\
- (\partial_t - \Delta) ([\Delta, \mathbbm{1}_\Theta \eta_{01}] z_\xi)(t,x).
\end{multline*}
This control function $h$ is localized in $(0,T) \times \omega$ due to the conditions on the support of $\eta_{01}$, $\eta_{12}$, $\eta_{23}$. This concludes the proof of Theorem \ref{Thm-insensitizing-Boundary}.

\begin{thebibliography}{99}
\bibitem{Alabau-Boussouira} F. Alabau-Boussouira - “Insensitizing exact controls for the scalar wave equation and exact controllability of 2-coupled cascade systems of PDE’s by a single control” , Math. Control Signals Systems 26 (2014) no. 1, p. 1-46
\bibitem{mousquetaires} F. Ammar Khodja, A. Benabdallah, M. González-Burgos \& L. de Teresa - “New phenomena for the null controllability of parabolic systems: minimal time and geometrical dependence” , J. Math. Anal. Appl. 444 (2016) no. 2, p. 1071-1113
\bibitem{Beauchard-Marbach-2017} K. Beauchard \& F. Marbach - “Quadratic obstructions to small-time local controllability for scalar-input systems” , J. Differential Equations 264 (2018) no. 5, p. 3704-3774
\bibitem{Bodart-GB} O. Bodart, M. González-Burgos \& R. Pérez-García - “Insensitizing controls for a heat equation with a nonlinear term involving the state and the gradient” , Nonlinear Anal. 57 (2004) no. 5-6, p. 687-711
\bibitem{Bodart-Fabre} O. Bodart \& C. Fabre - “Controls insensitizing the norm of the solution of a semilinear heat equation” , J. Math. Anal. Appl. 195 (1995) no. 3, p. 658-683
\bibitem{Boyer-HS-dT} F. Boyer, V. Hernández-Santamaría \& L. de Teresa - “Insensitizing controls for a semilinear parabolic equation: a numerical approach” , Math. Control Relat. Fields 9 (2019) no. 1, p. 117-158
\bibitem{Brezis} H. Brezis - Functional analysis, Sobolev spaces and partial differential equations , Universitext , Springer, New York, 2011
\bibitem{Carreno} N. Carreño - “Insensitizing controls for the Boussinesq system with no control on the temperature equation” , Adv. Differential Equations 22 (2017) no. 3-4, p. 235-258
\bibitem{Carreno-Guer-Gueye} N. Carreño, S. Guerrero \& M. Gueye - “Insensitizing controls with two vanishing components for the three-dimensional Boussinesq system” , ESAIM Control Optim. Calc. Var. 21 (2015) no. 1, p. 73-100
\bibitem{Carreno-Gueye} N. Carreño \& M. Gueye - “Insensitizing controls with one vanishing component for the Navier-Stokes system” , J. Math. Pures Appl. (9) 101 (2014) no. 1, p. 27-53
\bibitem{Chowdhury-Erv-2019} S. Chowdhury \& S. Ervedoza - “Open loop stabilization of incompressible Navier-Stokes equations in a 2d channel using power series expansion” , J. Math. Pures Appl. (9) 130 (2019), p. 301-346
\bibitem{CNL} J.-M. Coron - Control and nonlinearity , Math. Surveys and Monographs , vol. 136 , American Mathematical Society, Providence, RI, 2007
\bibitem{Delfour-Zolesio} M. C. Delfour \& J.-P. Zolésio - Shapes and geometries , Advances in Design and Control , vol. 22 , Society for Industrial and Applied Mathematics (SIAM), Philadelphia, PA, 2011
\bibitem{Ervedoza-Note-2020} S. Ervedoza - “Control issues and linear projection constraints on the control and on the controlled trajectory” , North-West. Eur. J. Math. 6 (2020), p. 165-197
\bibitem{Evans} L. C. Evans - Partial differential equations , Graduate Studies in Math. , vol. 19 , American Mathematical Society, Providence, RI, 2010
\bibitem{FernandezCaraZuazua1} E. Fernández-Cara \& E. Zuazua - “The cost of approximate controllability for heat equations: the linear case” , Adv. Differential Equations 5 (2000) no. 4-6, p. 465-514
\bibitem{FursikovImanuvilov} A. V. Fursikov \& O. Y. Imanuvilov - Controllability of evolution equations , Lect. Notes Series , vol. 34 , Seoul National University, Research Institute of Mathematics, Global Analysis Research Center, Seoul, 1996
\bibitem{Guerrero} S. Guerrero - “Controllability of systems of Stokes equations with one control force: existence of insensitizing controls” , Ann. Inst. H. Poincaré Anal. Non Linéaire 24 (2007) no. 6, p. 1029-1054
\bibitem{Gueye} M. Gueye - “Insensitizing controls for the Navier-Stokes equations” , Ann. Inst. H. Poincaré Anal. Non Linéaire 30 (2013) no. 5, p. 825-844
\bibitem{Henrot-Pierre-2005} A. Henrot \& M. Pierre - Shape variation and optimization , EMS Tracts in Math. , vol. 28 , European Mathematical Society, Zürich, 2018
\bibitem{Hormander-2003} L. Hörmander - The analysis of linear partial differential operators. I , Classics in Math. , Springer-Verlag, Berlin, 2003
\bibitem{Kavian-De-Teresa} O. Kavian \& L. de Teresa - “Unique continuation principle for systems of parabolic equations” , ESAIM Control Optim. Calc. Var. 16 (2010) no. 2, p. 247-274
\bibitem{Lions-1990} J.-L. Lions - “Quelques notions dans l’analyse et le contrôle de systèmes à données incomplètes” , in Proc. of the XIth Congress on Differential Equations and Appl./First Congress on Applied Math. (Málaga, 1989) , Univ. Málaga, Málaga, 1990, p. 43-54
\bibitem{Lions-1992} J.-L. Lions - “Remarks on approximate controllability” , J. Analyse Math. 59 (1992), p. 103-116
\bibitem{Lions-Sentinelles-1992} J.-L. Lions - Sentinelles pour les systèmes distribués à données incomplètes , Recherches en Math. Appl. , vol. 21 , Masson, Paris, 1992
\bibitem{LionsMagenes} J.-L. Lions \& E. Magenes - Problèmes aux limites non homogènes et applications. Vol. 1 , Travaux et Recherches Math. , vol. 17 , Dunod, Paris, 1968
\bibitem{LionsMagenes2} J.-L. Lions \& E. Magenes - Problèmes aux limites non homogènes et applications. Vol. 2 , Travaux et Recherches Math. , vol. 18 , Dunod, Paris, 1968
\bibitem{Lissy-Privat-Simpore} P. Lissy, Y. Privat \& Y. Simporé - “Insensitizing control for linear and semi-linear heat equations with partially unknown domain” , ESAIM Control Optim. Calc. Var. 25 (2019), article ID 50, 21 pages
\bibitem{Micu-Ortega-deTeresa} S. Micu, J. H. Ortega \& L. de Teresa - “An example of $\varepsilon $-insensitizing controls for the heat equation with no intersecting observation and control regions” , Appl. Math. Lett. 17 (2004) no. 8, p. 927-932
\bibitem{Tebou2} L. Tebou - “Locally distributed desensitizing controls for the wave equation” , Comptes Rendus Mathématique 346 (2008) no. 7-8, p. 407-412
\bibitem{Tebou1} L. Tebou - “Some results on the controllability of coupled semilinear wave equations: the desensitizing control case” , SIAM J. Control Optim. 49 (2011) no. 3, p. 1221-1238
\bibitem{deTeresa-00} L. de Teresa - “Insensitizing controls for a semilinear heat equation” , Comm. Partial Differential Equations 25 (2000) no. 1-2, p. 39-72
\bibitem{dT-Zuazua-09} L. de Teresa \& E. Zuazua - “Identification of the class of initial data for the insensitizing control of the heat equation” , Commun. Pure Appl. Anal. 8 (2009) no. 1, p. 457-471
\bibitem{TWBook} M. Tucsnak \& G. Weiss - Observation and control for operator semigroups , Birkhäuser Advanced Texts , Birkhäuser Verlag, Basel, 2009
\bibitem{Zuazua-97} E. Zuazua - “Finite-dimensional null controllability for the semilinear heat equation” , J. Math. Pures Appl. (9) 76 (1997) no. 3, p. 237-264
\end{thebibliography}
\end{document}
